% GENERATED FILE. Edit canonical lessons, then run npm run generate:books.
% canonical-source-hash: fnv1a64:7a6d1193e322962f
% canonical-lessons: HI-C113-kuchnahin, HI-C113-abhinahin, HI-C113-kafi, HI-C113-vahi, HI-C113-hosakta

\chapter{Nothing, Not Yet, Maybe}
\label{ch:hi-nothing-not-yet}
\hlchaptermodality{113}

\begin{chapteropening}
\textbf{By the end of this chapter:} \emph{I can answer with nothing, not yet, that's enough, that's the one, and maybe.}

Complete HI-C113-ho-sakta-hai: five short answers that close a question.
\end{chapteropening}

\section[\texorpdfstring{\hi{कुछ} \hi{नहीं} (kuch nahīṁ)}{कुछ नहीं (kuch nahīṁ)}]{\hi{कुछ} \hi{नहीं} (kuch nahīṁ) — nothing}

\label{lesson:HI-C113-kuchnahin}



\begin{quote}
Before the new one: say the Hindi for really, then the Hindi for oh.
\end{quote}

\subsection*{You'll want to know: \hi{कुछ} \hi{नहीं}}
\textbf{\hi{कुछ} \hi{नहीं}} (\emph{kuch nahīṁ}) — \textquotedblleft{}nothing\textquotedblright{}.

\textbf{\hi{क्या} \hi{हुआ}? — \hi{कुछ} \hi{नहीं}.} — \textquotedblleft{}what happened? — nothing\textquotedblright{}. It ends on the \textbf{\hi{नहीं}} you know.

\begin{cousinweb}
\textbf{\hi{कुछ} \hi{नहीं}} (\emph{kuch nahīṁ}) is two tadbhavas. \emph{kuch}, \textquotedblleft{}something\textquotedblright{}, comes from Sanskrit \emph{kiṁcid}, \textquotedblleft{}anything\textquotedblright{}, taking its \emph{-ch-} from the masculine \emph{kaścid} (CDIAL 3144); \emph{nahīṁ} comes from \emph{nahi}, \textquotedblleft{}surely not\textquotedblright{} (CDIAL 7035).
\end{cousinweb}

\begin{grammarlens}[title={What you've built}]
An answer that closes a question kindly.
\end{grammarlens}

\subsection*{Your turn}
\begin{itemize}
\raggedright
  \item \emph{Say it:} \emph{kuch nahīṁ}
  \item \emph{Say it:} \emph{kuch nahīṁ}, once more
  \item \emph{Say it:} say \emph{kuch nahīṁ}
  \item \emph{Recall:} say the Hindi for really, then the Hindi for oh, then say \emph{kuch nahīṁ} again
\end{itemize}

\begin{tcolorbox}[breakable,colback=teal!4,colframe=teal!35!black,title={Before you move on}]
What does \hi{कुछ} \hi{नहीं} mean? (\textquotedblleft{}Nothing\textquotedblright{}.) Say it once more.
\end{tcolorbox}
\section[\texorpdfstring{\hi{अभी} \hi{नहीं} (abhī nahīṁ)}{अभी नहीं (abhī nahīṁ)}]{\hi{अभी} \hi{नहीं} (abhī nahīṁ) — not yet, not now}

\label{lesson:HI-C113-abhinahin}



\begin{quote}
Before the new one: say the Hindi for oh, then the Hindi for nothing.
\end{quote}

\subsection*{You'll want to know: \hi{अभी} \hi{नहीं}}
\textbf{\hi{अभी} \hi{नहीं}} (\emph{abhī nahīṁ}) — \textquotedblleft{}not yet\textquotedblright{}, \textquotedblleft{}not now\textquotedblright{}.

Two words you own: \textbf{\hi{अभी}}, \textquotedblleft{}right now\textquotedblright{}, and \textbf{\hi{नहीं}}.

\begin{cousinweb}
\textbf{\hi{अभी} \hi{नहीं}} (\emph{abhī nahīṁ}) is \emph{abhī}, \textquotedblleft{}right now\textquotedblright{}, with \emph{nahīṁ}, \textquotedblleft{}not\textquotedblright{}. CDIAL lists \emph{ab} and \emph{abhī} under Sanskrit \emph{evam eva}, \textquotedblleft{}exactly so\textquotedblright{}, crossed with the pronoun stem \emph{a-} (CDIAL 2528); \emph{nahīṁ} comes from \emph{nahi}, \textquotedblleft{}surely not\textquotedblright{} (CDIAL 7035).
\end{cousinweb}

\begin{grammarlens}[title={What you've built}]
An answer from two old words.
\end{grammarlens}

\subsection*{Your turn}
\begin{itemize}
\raggedright
  \item \emph{Say it:} \emph{abhī nahīṁ}
  \item \emph{Say it:} \emph{abhī nahīṁ}, once more
  \item \emph{Say it:} say \emph{abhī nahīṁ}
  \item \emph{Recall:} say the Hindi for oh, then the Hindi for nothing, then say \emph{abhī nahīṁ} again
\end{itemize}

\begin{tcolorbox}[breakable,colback=teal!4,colframe=teal!35!black,title={Before you move on}]
What does \hi{अभी} \hi{नहीं} mean? (\textquotedblleft{}Not yet, not now\textquotedblright{}.) Say it once more.
\end{tcolorbox}
\section[\texorpdfstring{\hi{काफ़ी} (kāfī)}{काफ़ी (kāfī)}]{\hi{काफ़ी} (kāfī) — enough; quite}

\label{lesson:HI-C113-kafi}



\begin{quote}
Before the new one: say the Hindi for nothing, then the Hindi for not yet.
\end{quote}

\subsection*{You'll want to know: \hi{काफ़ी}}
\textbf{\hi{काफ़ी}} (\emph{kāfī}) — \textquotedblleft{}enough\textquotedblright{}, and before an adjective, \textquotedblleft{}quite\textquotedblright{}.

\textbf{\hi{बस}, \hi{काफ़ी} \hi{है}} — \textquotedblleft{}that's enough\textquotedblright{}, with the \textbf{\hi{बस}} you know. \textbf{\hi{काफ़ी} \hi{अच्छा}} — \textquotedblleft{}quite good\textquotedblright{}.

\begin{cousinweb}
\textbf{\hi{काफ़ी}} (\emph{kāfī}) is a loan from Arabic \emph{kāfī}, \textquotedblleft{}sufficient\textquotedblright{}, which reached Hindi through Persian. Its Arabic root is \emph{k-f-y} (Lane).
\end{cousinweb}

\begin{grammarlens}[title={What you've built}]
Enough, and quite.
\end{grammarlens}

\subsection*{Your turn}
\begin{itemize}
\raggedright
  \item \emph{Say it:} \emph{kāfī}
  \item \emph{Say it:} \emph{kāfī}, once more
  \item \emph{Say it:} say \emph{bas, kāfī hai}
  \item \emph{Recall:} say the Hindi for nothing, then the Hindi for not yet, then say \emph{kāfī} again
\end{itemize}

\begin{tcolorbox}[breakable,colback=teal!4,colframe=teal!35!black,title={Before you move on}]
What does \hi{काफ़ी} mean? (\textquotedblleft{}Enough; quite\textquotedblright{}.) Say it once more.
\end{tcolorbox}
\section[\texorpdfstring{\hi{वही} (vahī)}{वही (vahī)}]{\hi{वही} (vahī) — that very one, the same}

\label{lesson:HI-C113-vahi}



\begin{quote}
Before the new one: say the Hindi for not yet, then the Hindi for enough.
\end{quote}

\subsection*{You'll want to know: \hi{वही}}
\textbf{\hi{वही}} (\emph{vahī}) — \textquotedblleft{}that very one\textquotedblright{}, \textquotedblleft{}the same\textquotedblright{}.

It is \textbf{\hi{वह}}, \textquotedblleft{}that\textquotedblright{}, with \textbf{\hi{ही}}, \textquotedblleft{}exactly\textquotedblright{}. When someone names what you meant: \textbf{\hi{हाँ}, \hi{वही}!}

\begin{cousinweb}
\textbf{\hi{वही}} (\emph{vahī}) is \textbf{\hi{वह}}, \textquotedblleft{}that\textquotedblright{}, with the emphatic \emph{-hī}: \textquotedblleft{}that one and no other\textquotedblright{}. CDIAL lists \emph{vah} under Sanskrit \emph{asau}, \textquotedblleft{}that\textquotedblright{} (CDIAL 972).
\end{cousinweb}

\begin{grammarlens}[title={What you've built}]
That one, exactly.
\end{grammarlens}

\subsection*{Your turn}
\begin{itemize}
\raggedright
  \item \emph{Say it:} \emph{vahī}
  \item \emph{Say it:} \emph{vahī}, once more
  \item \emph{Say it:} say \emph{hā̃, vahī!}
  \item \emph{Recall:} say the Hindi for not yet, then the Hindi for enough, then say \emph{vahī} again
\end{itemize}

\begin{tcolorbox}[breakable,colback=teal!4,colframe=teal!35!black,title={Before you move on}]
What does \hi{वही} mean? (\textquotedblleft{}That very one, the same\textquotedblright{}.) Say it once more.
\end{tcolorbox}
\section[\texorpdfstring{\hi{हो} \hi{सकता} \hi{है} (ho saktā hai)}{हो सकता है (ho saktā hai)}]{\hi{हो} \hi{सकता} \hi{है} (ho saktā hai) — maybe, it could be}

\label{lesson:HI-C113-hosakta}



\begin{quote}
Before the new one: say the Hindi for enough, then the Hindi for that very one.
\end{quote}

\subsection*{You'll want to know: \hi{हो} \hi{सकता} \hi{है}}
\textbf{\hi{हो} \hi{सकता} \hi{है}} (\emph{ho saktā hai}) — \textquotedblleft{}maybe\textquotedblright{}, literally \textquotedblleft{}it can be\textquotedblright{}.

It is \textbf{\hi{होना}}, \textquotedblleft{}to be\textquotedblright{}, and the \textbf{\hi{सकना}} of \textquotedblleft{}can\textquotedblright{}. Compare \textbf{\hi{शायद}}, \textquotedblleft{}perhaps\textquotedblright{}.

\begin{cousinweb}
\textbf{\hi{हो} \hi{सकता} \hi{है}} (\emph{ho saktā hai}) is \textquotedblleft{}it can be\textquotedblright{}. \emph{ho} is a form of \emph{honā}, \textquotedblleft{}to be\textquotedblright{}, from Sanskrit \emph{bhavati} (CDIAL 9416), and \emph{saknā}, \textquotedblleft{}to be able\textquotedblright{}, comes from \emph{śaknoti}, \textquotedblleft{}is able\textquotedblright{} (CDIAL 12252).
\end{cousinweb}

\begin{grammarlens}[title={What you've built}]
Nothing, not yet, enough, that one, maybe.
\end{grammarlens}

\subsection*{Your turn}
\begin{itemize}
\raggedright
  \item \emph{Say it:} \emph{ho saktā hai}
  \item \emph{Say it:} \emph{ho saktā hai}, once more
  \item \emph{Say it:} say \emph{ho saktā hai}
  \item \emph{Recall:} say the Hindi for enough, then the Hindi for that very one, then say \emph{ho saktā hai} again
\end{itemize}

\begin{tcolorbox}[breakable,colback=teal!4,colframe=teal!35!black,title={Before you move on}]
What does \hi{हो} \hi{सकता} \hi{है} mean? (\textquotedblleft{}Maybe, it could be\textquotedblright{}.) Say it once more.
\end{tcolorbox}
